\subsection{Explicit protection against reference loss}
We keep the early proposal generator fixed and train the probability
that an early action replaces the only existing easy reference without
adding one. A factored score multiplies the predicted probability that
the original support is unique by the conditional probability of no new
reference. Both see the current candidate and observed prefix. They
retain proposals below the 75th percentile of risk scores on donor
proposals. A rejection leaves the original policy active. This percentile
sets a development budget, not a calibrated risk bound.

\begin{center}
\begin{tabular}{lc}
Arm & Edits / joint / undefined \\
\hline
Original & \RefLossOriginalCounts \\
Frozen early policy & \RefLossFrozenCounts \\
Direct loss guard & \RefLossDirectCounts \\
Factored guard & \RefLossFactoredCounts \\
Easy rank, direct-matched & \RefLossMatchedCounts \\
Exact count, hindsight & \RefLossExactCounts \\
No loss, up to count, hindsight & \RefLossNoLossCounts \\
\end{tabular}
\end{center}

The direct guard keeps \RefLossAccepted{} of \RefLossProposed{}
proposals. Reference losses decline from \RefLossBefore{} to
\RefLossAfter{}, but the conservative utility difference versus the
frozen early policy is \RefLossFrozenEffect{}. Against easy ranking
matched separately on additions and replacements in each location
and seed, it is \RefLossMatchedEffect{}. Intervals are paired by
location. Even hindsight cannot fill \RefLossShortfall{} of the
matched slots without losing a reference. With smaller budgets,
the no-loss hindsight comparison adds only four joint passes
over the original. It is an outcome-aware diagnostic, not a causal
improvement. The 2\% recording rule is unchanged.
